\section{Conclusion}\label{sec:conclusion}

The question behind this thesis was whether the derivative-based theory of regular
languages still works once exact symbol matching is dropped and symbols are allowed
to match by similarity instead, and whether the core of that theory can be checked
by a machine rather than trusted on paper. The work in the preceding sections
answers yes to both.

On the crisp side the thesis gives a self-contained account of derivatives of
regular languages and of regular expressions. The result everything else rests on is
Theorem~\ref{the:two-derivative-mapping}, which says that the syntactic derivative
of an expression and the semantic derivative of its language describe the same
thing once an expression is read as the language it denotes. From there it is a
short step to the fact that derivatives stay regular and that an expression has only
finitely many derivative languages (Theorem~\ref{the:derivative-set-is-finite}).
That finiteness is exactly what makes the derivative automaton possible: taking
canonical derivatives as states and differentiation as the transition yields a
deterministic automaton that recognises the right language
(Theorem~\ref{the:dfa-correctness}) and that a terminating algorithm actually
builds.

The fuzzy layer is where the thesis makes its own contribution. Once equality of
symbols is replaced by a fuzzy similarity relation together with a cut value $\mu$,
fuzzy derivatives of both languages and restricted regular expressions arise
(Definitions~\ref{def:fuzzy-deriv-lang} and~\ref{def:fuzzy-deriv-expr}). The idea
that makes the whole layer tractable is small but decisive: under the minimum
T-norm, cut-similarity turns out to be an equivalence relation
(Proposition~\ref{prop:cut-equivalence}). Once that is in hand, a fuzzy derivative
over $\Sigma$ is just an ordinary derivative over the quotient alphabet
$\Sigma^{\mu}$ obtained by collapsing $\mu$-similar symbols into a single class
(Theorem~\ref{the:quotient-fuzzy-derivative}). This quotient picture is what lets
the crisp theory carry over almost without extra effort, so regularity
(Theorem~\ref{the:fuzzy-regular}), the correspondence between syntax and semantics
(Theorem~\ref{the:fuzzy-correspondence}), finiteness
(Theorem~\ref{the:fuzzy-finite}) and fuzzy DFA correctness
(Theorem~\ref{the:fuzzy-dfa}) all follow. In plain terms, the fuzzy automaton
accepts precisely the words that come within tolerance $\mu$ of some word in the
original language.

This quotient construction is really the point of the thesis. It does more than
confirm that the fuzzy theorems happen to hold; it explains why fuzzification by
cut-similarity behaves so well in the first place, and it makes the fuzzy derivative
something computable, since the computation collapses to ordinary derivatives
over a finite quotient alphabet. The verification side carries its own weight.
Mechanising the crisp theory in Rocq turns the central claims into statements a
proof assistant has checked, covering the syntax and semantics, nullability,
canonicalisation and its correctness, the derivative correspondence and DFA
correctness. That assurance matters most where it is hardest to be sure by hand,
which is in the constructive finiteness arguments.

The fuzzy expression theory is
developed on the restricted fragment, with concatenation, union and star, where
syntactic similarity behaves monotonically and the quotient argument goes through
cleanly (Remark~\ref{rem:syntax-sensitive}). It rests on the idempotence of the
minimum T-norm, which is what turns cut-similarity into an equivalence and lets the
quotient construction do its work, and alongside it the thesis gives a definition of
the fuzzy derivative that already reads for a general T-norm. The crisp layer is
mechanised in Rocq end to end, and the similarity infrastructure that the fuzzy
results build on is formalised there as well.

These same boundaries mark the natural next steps. The most immediate is to carry the
fuzzy layer into Rocq on top of the similarity infrastructure that is already in
place, and to settle the targeted Kuratowski-finiteness result for canonical
derivatives. Beyond that, it would be interesting to study fuzzy derivatives for
general, non-idempotent left-continuous T-norms, where cut relations stop being
equivalences and the clean quotient construction has to give way to something more
subtle. A further direction is to connect this similarity-based account more
carefully to the membership-weight tradition of fuzzy
automata~\cite{LP06,SC12,GJT17,SilvaBonchiBonsangueRutten2011} and to the
quantitative pattern-matching line of
work~\cite{AlvesKesnerRamos2025}, and to develop the fuzzy derivative automaton as a
genuine fuzzy transition system. Taken together, these steps would turn what is
developed here into a fully mechanised and more widely applicable framework for
approximate matching.
